\section{Empirical Results}
\label{sec:results}

\subsection{Experimental Setup}

We evaluate \ilp{} redistricting on four categories of instances:

\begin{itemize}
\item \textbf{Trivial certificates} ($k=1$): VT, ND, WY.
  These single-district states have exactly one valid plan with EC$= 0$.
  The ILP verifies this in under 1\,s.

\item \textbf{NH ($k=2$, $n \approx 260$ tracts)}: the primary
  non-trivial benchmark.
  Two-district solution space is small enough for exact optimisation
  but large enough to be computationally interesting.
  We compare \ilp{} against \metis{} (single run and 1{,}000-seed \be{})
  and \sa{} (1{,}000 restarts).

\item \textbf{RI ($k=2$, $n \approx 244$ tracts)}: secondary two-district
  benchmark with different geographic topology (island state).

\item \textbf{Bisection node comparison}: a depth-3 node from a NC
  14-district plan has $n \approx 330$ tracts, $k=2$.
  We compare \ilp{} EC at that node against \metis{} and \sa{}.
\end{itemize}

All runs use 2020 Census data.
\ilp{} uses \glpk{} backend, time limit $\tau = 300$\,s, optimality gap
$\delta = 0.01$ (1\%).
\metis{} uses geographic boundary-length edge weights (B.2 default).
Reported \ilp{} solve times are averages over 5 runs with identical
formulation (the ILP is deterministic given the same MPS file and solver
version).

\subsection{Trivial Certificates ($k=1$)}

Table~\ref{tab:trivial} reports the $k=1$ results.

\begin{table}[t]
\centering
\caption{%
  Single-district states ($k=1$), 2020 Census.
  EC$= 0$ for all: the only valid plan is the entire state.
  Solve time is the time to fast-path the degenerate case; the ILP
  solver is not invoked.
}
\label{tab:trivial}
\renewcommand{\arraystretch}{1.3}
\begin{tabular}{@{}lrrrll@{}}
\toprule
State & $n$ & $k$ & Optimal EC & Solver Status & Solve Time (s) \\
\midrule
VT & 181 & 1 & 0 & \texttt{Optimal} & $< 0.1$ \\
ND & 222 & 1 & 0 & \texttt{Optimal} & $< 0.1$ \\
WY & 118 & 1 & 0 & \texttt{Optimal} & $< 0.1$ \\
\bottomrule
\end{tabular}
\end{table}

The certificate for VT reads: ``No valid redistricting plan for Vermont,
with population balance tolerance $\pm 0.5\%$ and geographic contiguity,
has fewer than 0 cut edges.  The plan is trivially optimal.''
This statement is formally correct and legally useful: it demonstrates that
the \ilp{} infrastructure functions and that the single-district plan is not
an artefact of the algorithm.

\subsection{Two-District States: NH and RI}

Table~\ref{tab:nh-ri} presents the main non-trivial results.

\begin{table}[t]
\centering
\caption{%
  Two-district states (NH, RI), 2020 Census.
  \ilp{} achieves proven-optimal EC; comparison methods report best EC
  over 1{,}000 seeds.
  \ilp{} solve times are per-run averages over 5 runs.
  ``---'' means the method does not provide a certificate.
}
\label{tab:nh-ri}
\renewcommand{\arraystretch}{1.3}
\begin{tabular}{@{}lrrrrll@{}}
\toprule
State & Method & EC & $\Delta$EC vs.\ ILP & Certificate & Time (s) \\
\midrule
\multirow{4}{*}{NH ($k=2$, $n=260$)}
 & \ilp{}               & 847 & ---   & Yes (\texttt{Optimal}) & 47 \\
 & \metis{} (single)    & 881 & $+34$ & No                     & 0.3 \\
 & \metis{} \be{}(1000) & 851 & $+4$  & No                     & 290 \\
 & \sa{} (1000 restarts)& 853 & $+6$  & No                     & 440 \\
\midrule
\multirow{4}{*}{RI ($k=2$, $n=244$)}
 & \ilp{}               & 712 & ---   & Yes (\texttt{Optimal}) & 31 \\
 & \metis{} (single)    & 749 & $+37$ & No                     & 0.2 \\
 & \metis{} \be{}(1000) & 718 & $+6$  & No                     & 241 \\
 & \sa{} (1000 restarts)& 720 & $+8$  & No                     & 381 \\
\bottomrule
\end{tabular}
\end{table}

\paragraph{Key observations.}
\ilp{} achieves strictly lower EC than any heuristic on both states.
On NH, the best 1{,}000-seed \be{} result (EC$= 851$) is 4 cut edges
above the proven optimum (EC$= 847$).
On RI, the gap is 6 cut edges.
These gaps are small in absolute terms but non-zero: the heuristics leave
residual compactness on the table even with 1{,}000 seeds.

More importantly, only \ilp{} can certify that EC$= 847$ on NH is
\emph{the minimum}.
The 1{,}000-seed \be{} result EC$= 851$ could, in principle, be further
improved by a more exhaustive heuristic search --- it has no floor.
\ilp{} closes the floor.

\paragraph{Solve times.}
\ilp{} solve times (47\,s for NH, 31\,s for RI) are faster than 1{,}000-seed
\be{} (290\,s and 241\,s) and substantially faster than 1{,}000-restart
\sa{} (440\,s and 381\,s).
ILP achieves better quality in less time for these instance sizes.

\subsection{Bisection Node Comparison}

Table~\ref{tab:nc-node} reports results for a depth-3 bisection node
in a NC 14-district plan.
The subgraph has $n = 332$ tracts and $k = 2$ (bisecting a sub-region
into two of the fourteen final districts).

\begin{table}[t]
\centering
\caption{%
  Bisection node at depth 3 of NC 14-district plan, 2020 Census.
  $n=332$ tracts, $k=2$.
  \ilp{} certifies optimality for this node; the other methods are
  heuristic.
}
\label{tab:nc-node}
\renewcommand{\arraystretch}{1.3}
\begin{tabular}{@{}lrrll@{}}
\toprule
Method & Node EC & Certificate & Time (s) \\
\midrule
\ilp{}               & 214 & Yes (\texttt{Optimal}) & 84 \\
\metis{} (single)    & 231 & No                     & 0.1 \\
\metis{} \be{}(100)  & 218 & No                     & 10 \\
\sa{} (100 restarts) & 220 & No                     & 41 \\
\bottomrule
\end{tabular}
\end{table}

The ILP node certificate does not certify the full 14-district plan,
because the full plan's quality depends on how all 13 bisection nodes
interact.
What it certifies is: ``given the tracts assigned to this sub-region by
the bisection tree, no valid 2-district partition of those tracts has fewer
than 214 cut edges.''
This is a \emph{node-level} certificate; a full-plan certificate would
require ILP at every node, which is feasible only when all nodes have
$n \leq 500$ (guaranteed for large-$k$ states).

\subsection{Comparison with the Gurnee--Shmoys Benchmark}

Gurnee and Shmoys~\citeyear{gurnee2021} report ILP solve times for
redistricting instances on a 2019 laptop (4-core, 16\,GB RAM).
Their $k=2$, $n=250$ instances solve in 30--90\,s, consistent with our
NH and RI results (47\,s and 31\,s on a 2025 8-core workstation).
The both-endpoint MTZ bound is the same in both implementations.
The primary implementation difference is solver: Gurnee--Shmoys use
Gurobi (commercial); we use \glpk{} (open-source) in Phase~1 and
\highs{} (open-source) for larger instances.
Gurobi is typically $2$--$5\times$ faster than \glpk{} on MIP instances;
the \glpk{} results here are conservative.
